\chapter{Introduction}
\label{chap:intro}


Courant and Robbins, in \href{https://en.wikipedia.org/wiki/What_Is_Mathematics\%3F}{What is Mathematics? (1941)}, present mathematics not as a dry collection of formulas and tools, but as a living, creative discipline rooted in human thought and curiosity. For them, mathematics is both a pathway for understanding the natural world and an autonomous intellectual pursuit that reveals structures of order, beauty, and generality. They stress that its essence lies in the interplay between abstraction and concrete problem-solving: starting from simple, practical problems, mathematics ascends to general concepts and theories that then illuminate new domains.

They emphasize accessibility and unity: mathematics belongs to everyone who is willing to think rigorously, and its spirit combines logic with imagination. Rather than reducing it to calculation or technical skill, Courant and Robbins describe mathematics as “an expression of the human mind” where precision, creativity, and aesthetic appreciation converge. Their central idea is that mathematics is at once useful, philosophical, and artistic — simultaneously a language of science, a training ground for reasoning, and a source of intellectual delight.

Early mathematics was computational when the emphasis was on finding methods to obtain solutions. However, over the years, the disciplines of mathematics and computer science -- the subject of designing algorithms for problem solving -- have diverged. In mathematics abstraction is symbolic and logical. It seeks general structures, patterns, and proofs independent of implementation. It often endeavours to seek and  capture  common structures across different abstractions. The primary aim is truth and understanding -- developing rigorous proofs, ensuring logical consistency, and uncovering general laws. Utility often follows from this pursuit but is not always the main driver. In contrast, the role of abstraction in computational thinking  is more operational and algorithmic. It emphasizes creating computational process models for natural, social and even abstract phenomena for operational analysis. The primary aim is to construct effective procedures -- designing algorithms that solve problems efficiently, often under constraints of time, memory, and real-world complexity. The power lies in execution and exploration — trying to  reveal insights about systems too complex to solve analytically.  Both have become fundamental strands of epistemology that are essential for critical scientific thinking.

Data-driven inference represents a third way of knowing, distinct from the deductive rigour of mathematics and the constructive procedures of computational thinking. As practiced in modern data science and machine learning, it seeks knowledge not by proving theorems or designing explicit algorithms, but by discovering patterns and regularities directly from empirical data. Its epistemic core is induction at scale: hypotheses, models, or predictors are justified by their ability to capture hidden correlations and to generalize to new observations. Unlike mathematics, correctness is not absolute, and unlike computational thinking, procedures are not always fully transparent. Instead, credibility arises from empirical adequacy — the degree to which models explain, predict, or align with observed phenomena. This mode of inference expands our epistemic toolkit for a world where complexity and abundance of data overwhelm deductive or constructive methods, but it also brings new philosophical challenges: uncertainty about correctness, bias, and the gap between correlation and causation.

In this course we will try to cover some fundamentals of all of the above.

